% =============================================================================
% ANEXO G — Especificación detallada de endpoints (payloads y validaciones)
% Extraído de los DTO reales (records con Jakarta Bean Validation) de ethoshubB.
% =============================================================================
\chapter{Anexo G — Especificación Detallada de Endpoints}
\label{cap:anexo_ep_detalle}

Este anexo documenta, dominio por dominio, los \textit{payloads} de petición y respuesta y las
reglas de validación reales de la API \comando{ethoshubB}, extraídas de los \comando{record} de
DTO anotados con Jakarta Bean Validation. Todas las respuestas se envuelven en el tipo genérico
\comando{ApiResponse<T>} (Sección~\ref{sec:anexo_ep_envelope}).\par

\section{Envoltorio de Respuesta \texttt{ApiResponse}}
\label{sec:anexo_ep_envelope}

Toda respuesta de la API comparte la misma estructura: un indicador de éxito, el código de
estado reflejado en el cuerpo, un mensaje legible, el \textit{payload} (\comando{data}, nulo en
error) y la lista de errores (\comando{errors}, nula en éxito).

\begin{lstlisting}[language=json, caption={Estructura común \comando{ApiResponse<T>}.}, label={lst:anexo_apiresponse}]
{ "success": true,
  "status": 200,
  "message": "Login successful",
  "data": { },
  "errors": null }
\end{lstlisting}

\section{Autenticación (\texttt{/api/auth})}
\label{sec:anexo_ep_auth_detalle}

El \comando{AuthController} expone el registro, el inicio de sesión y la sincronización OAuth. La
tabla~\ref{tab:anexo_auth_val} detalla las validaciones de los DTO de entrada.

\begin{lstlisting}[language=json, caption={\comando{POST /api/auth/register} --- petición y respuesta.}, label={lst:anexo_register}]
// Request (RegisterRequest)
{ "email": "user@example.com", "password": "S3cur3P@ss!", "role": "PROFESSIONAL" }
// Response (ApiResponse<AuthResponse>)
{ "success": true, "status": 201, "message": "Registered",
  "data": { "token": "eyJhbGciOiJIUzI1...", "profileId": "550e8400-...", "email": "user@example.com" },
  "errors": null }
\end{lstlisting}

\begin{table}[!ht]
    \centering
    \renewcommand{\arraystretch}{1.3}
    \fontsize{9}{10.8}\selectfont\sffamily
    \begin{tabularx}{\textwidth}{|l|l|X|}
        \hline
        \rowcolor{colorPrimario}
        \textcolor{white}{\textbf{Campo}} & \textcolor{white}{\textbf{Validación}} & \textcolor{white}{\textbf{Regla}} \\ \hline
        \comando{email} & \comando{@NotBlank @Email} & Requerido; formato de correo válido. \\ \hline
        \rowcolor{grisTecnico}
        \comando{password} & \comando{@NotBlank} & Requerido; mínimo 8 caracteres. \\ \hline
        \comando{role} & --- & \comando{PROFESSIONAL} o \comando{RECRUITER}. \\ \hline
    \end{tabularx}
    \caption{Validaciones de \comando{RegisterRequest} / \comando{LoginRequest}.}
    \label{tab:anexo_auth_val}
\end{table}

\section{Seguridad y OTP (\texttt{/api/v1/auth})}
\label{sec:anexo_ep_security_detalle}

El \comando{SecurityController} gestiona las operaciones sensibles mediante OTP. La
tabla~\ref{tab:anexo_security_ep} enumera sus endpoints reales.

\begin{table}[!ht]
    \centering
    \renewcommand{\arraystretch}{1.3}
    \fontsize{9}{10.8}\selectfont\sffamily
    \begin{tabularx}{\textwidth}{|l|l|X|}
        \hline
        \rowcolor{colorPrimario}
        \textcolor{white}{\textbf{Método}} & \textcolor{white}{\textbf{Ruta}} & \textcolor{white}{\textbf{Propósito}} \\ \hline
        POST & \comando{/request-password-change} & Solicita OTP para cambio de contraseña. \\ \hline
        \rowcolor{grisTecnico}
        POST & \comando{/request-email-change} & Solicita OTP para cambio de email. \\ \hline
        POST & \comando{/request-account-delete} & Solicita OTP para baja de cuenta. \\ \hline
        \rowcolor{grisTecnico}
        POST & \comando{/verify-otp} & Verifica el código OTP. \\ \hline
        GET & \comando{/confirm-email-change} & Confirma el cambio vía enlace. \\ \hline
        \rowcolor{grisTecnico}
        DELETE & \comando{/account} & Ejecuta \comando{delete\_profile\_cascade}. \\ \hline
        GET & \comando{/export-data} & Exporta los datos del usuario. \\ \hline
    \end{tabularx}
    \caption{Endpoints reales de \comando{SecurityController} (\comando{/api/v1/auth}).}
    \label{tab:anexo_security_ep}
\end{table}

\section{Habilidades (\texttt{/api/skills}, \texttt{/api/profiles/\{id\}/skills})}
\label{sec:anexo_ep_skills_detalle}

\begin{lstlisting}[language=json, caption={\comando{POST /api/profiles/\{id\}/skills/hard} y \comando{.../soft}.}, label={lst:anexo_skills}]
// HardSkillRequest
{ "tagId": "a1b2c3d4-...", "level": "Advanced" }
// SoftSkillRequest
{ "title": "Comunicación", "description": "Casos de éxito en equipos distribuidos." }
\end{lstlisting}

\begin{table}[!ht]
    \centering
    \renewcommand{\arraystretch}{1.3}
    \fontsize{9}{10.8}\selectfont\sffamily
    \begin{tabularx}{\textwidth}{|l|l|X|}
        \hline
        \rowcolor{colorPrimario}
        \textcolor{white}{\textbf{DTO}} & \textcolor{white}{\textbf{Campo}} & \textcolor{white}{\textbf{Validación}} \\ \hline
        HardSkillRequest & \comando{tagId} & \comando{@NotBlank}. \\ \hline
        \rowcolor{grisTecnico}
        HardSkillRequest & \comando{level} & \comando{@Size(max=50)}. \\ \hline
        SoftSkillRequest & \comando{title} & \comando{@NotBlank @Size(max=100)}. \\ \hline
        \rowcolor{grisTecnico}
        SoftSkillRequest & \comando{description} & \comando{@Size(max=500)}. \\ \hline
        SkillTagCreateRequest & \comando{name} / \comando{category} & \comando{@NotBlank @Size(max=100)} / \comando{@Size(max=100)}. \\ \hline
    \end{tabularx}
    \caption{Validaciones de los DTO de habilidades.}
    \label{tab:anexo_skills_val}
\end{table}

\section{Proyectos (\texttt{/api/projects})}
\label{sec:anexo_ep_projects_detalle}

\begin{lstlisting}[language=json, caption={\comando{POST /api/projects} --- \comando{ProjectRequest}.}, label={lst:anexo_project}]
{ "title": "Portafolio EthosHub",
  "description": "SPA + API monolítica para portafolios y talento.",
  "repositoryUrl": "https://github.com/bytebusters/ethoshub",
  "isFeatured": true, "visibility": "public", "category": "Web",
  "media": [ { "type": "image", "url": "https://...", "title": "Home" } ] }
\end{lstlisting}

\begin{table}[!ht]
    \centering
    \renewcommand{\arraystretch}{1.3}
    \fontsize{9}{10.8}\selectfont\sffamily
    \begin{tabularx}{\textwidth}{|l|l|X|}
        \hline
        \rowcolor{colorPrimario}
        \textcolor{white}{\textbf{Campo}} & \textcolor{white}{\textbf{Validación}} & \textcolor{white}{\textbf{Notas}} \\ \hline
        \comando{title} & \comando{@NotBlank @Size(max=200)} & Título del proyecto. \\ \hline
        \rowcolor{grisTecnico}
        \comando{description} & \comando{@Size(max=5000)} & Descripción extensa. \\ \hline
        \comando{repositoryUrl} & \comando{@Size(max=500)} & URL del repositorio. \\ \hline
        \rowcolor{grisTecnico}
        \comando{visibility} & \comando{@Size(max=50)} & \comando{public} / \comando{private}. \\ \hline
        \comando{category} & \comando{@Size(max=100)} & Categoría. \\ \hline
        \rowcolor{grisTecnico}
        \comando{media[]} & --- & Lista de \comando{MediaRequest} (\comando{type}, \comando{url}, \comando{title}). \\ \hline
    \end{tabularx}
    \caption{Validaciones de \comando{ProjectRequest}.}
    \label{tab:anexo_project_val}
\end{table}

\section{Experiencia y Conexiones}
\label{sec:anexo_ep_otros_detalle}

\begin{table}[!ht]
    \centering
    \renewcommand{\arraystretch}{1.3}
    \fontsize{9}{10.8}\selectfont\sffamily
    \begin{tabularx}{\textwidth}{|l|l|X|}
        \hline
        \rowcolor{colorPrimario}
        \textcolor{white}{\textbf{Método}} & \textcolor{white}{\textbf{Ruta}} & \textcolor{white}{\textbf{Propósito}} \\ \hline
        GET & \comando{/api/v1/work-experiences/profile/\{id\}} & Lista la experiencia del perfil. \\ \hline
        \rowcolor{grisTecnico}
        POST/PUT & \comando{/api/v1/work-experiences[/\{id\}]} & Crea/actualiza experiencia. \\ \hline
        DELETE & \comando{/api/v1/work-experiences/\{id\}/profile/\{id\}} & Baja lógica de experiencia. \\ \hline
        \rowcolor{grisTecnico}
        GET/POST & \comando{/api/v1/connections} & Lista/crea conexiones URL-first. \\ \hline
        POST & \comando{/api/v1/connections/\{id\}/disconnect|reconnect|sync} & Gestiona el estado de la conexión. \\ \hline
        \rowcolor{grisTecnico}
        DELETE & \comando{/api/v1/connections/\{id\}} & Elimina la conexión. \\ \hline
    \end{tabularx}
    \caption{Endpoints reales de experiencia y conexiones.}
    \label{tab:anexo_otros_ep}
\end{table}

\begin{NotaTecnica}
Los errores de validación de Jakarta se transforman en respuestas \comando{ApiResponse} con
\comando{success=false}, \comando{status=400} y la lista \comando{errors} poblada con los
mensajes de cada restricción incumplida, gestionados por el manejador global de excepciones de
\comando{shared/}.
\end{NotaTecnica}

\clearpage
